\documentclass[fleqn]{beamer}

\mode<presentation> {
\usetheme{Madrid}
\usecolortheme{dolphin}
}

\usepackage[utf8]{inputenc}
\usepackage[brazilian]{babel}
\usepackage{amsmath}
\usepackage{amssymb}
\usepackage{booktabs}
\usepackage{listings}
\usepackage{xcolor}
\usepackage{tikz}
\usetikzlibrary{arrows.meta}
\usepackage{pgfplots}
\pgfplotsset{compat=1.15}
\usepackage{graphicx}

\lstset{
	basicstyle=\ttfamily\scriptsize,
	breaklines=true,
	keywordstyle=\color{blue!60!black},
	commentstyle=\color{gray},
	stringstyle=\color{green!45!black},
	showstringspaces=false,
	columns=fullflexible,
	frame=single,
	framesep=4pt,
	aboveskip=6pt,
	belowskip=2pt,
}

% ---------- Figuras reutilizáveis ----------

% Paisagem unidimensional de um problema de maximização:
% três máximos locais (x = 1,8; 5,8; 9,0), sendo 5,8 o global.
\newcommand{\paisagem}{%
	\draw[->] (-0.2,0) -- (9.8,0) node[right] {\scriptsize $x \in S$};
	\draw[->] (0,-0.1) -- (0,4.2) node[above] {\scriptsize $f(x)$};
	\draw[thick,blue!60!black] plot[smooth] coordinates {%
		(0.3,1.0) (1.0,2.0) (1.8,2.6) (2.6,2.0) (3.6,0.8) (4.6,1.6)
		(5.8,3.6) (6.8,2.4) (7.6,1.8) (8.4,2.4) (9.0,2.6) (9.6,2.2)};
}

% Paisagem bidimensional (superfície 3D) com os mesmos três máximos:
% (2,2) e (8,2) locais, (6,6) global. Recebe o conteúdo extra do eixo.
\newcommand{\paisagemtres}[1]{%
	\begin{tikzpicture}
	\begin{axis}[width=0.98\linewidth,height=5.2cm,view={-30}{35},
		hide axis,colorbar,colorbar style={width=2mm,ytick=\empty,
			title={\scriptsize $f(x)$},title style={yshift=-1mm}},
		colormap={custo}{color=(red) color=(yellow) color=(green) color=(cyan) color=(blue) color=(violet)},
		domain=0:10,y domain=0:10,samples=26,zmin=0,zmax=4.5]
		\addplot3[surf,shader=interp]
			{2.6*exp(-((x-2)^2+(y-2)^2)/1.5)+3.6*exp(-((x-6)^2+(y-6)^2)/1.5)
			+2.6*exp(-((x-8)^2+(y-2)^2)/1.5)};
		#1
	\end{axis}
	\end{tikzpicture}%
}

% Linha de um vetor binário (#1 = coordenada y, #2 = lista de bits)
\tikzset{bit/.style={draw,minimum size=0.6cm,inner sep=0pt,font=\small},
	bit0/.style={bit,fill=white}, bit1/.style={bit,fill=green!25}}
\newcommand{\vetor}[2]{%
	\foreach \v [count=\i] in {#2}{%
		\node[bit\v] at (\i*0.65,#1) {\v};
	}%
}
% Contorno em destaque de uma posição do vetor (#1 = y, #2 = índice)
\newcommand{\destaque}[2]{%
	\draw[red,very thick] ({#2*0.65-0.3},{#1-0.3}) rectangle ({#2*0.65+0.3},{#1+0.3});
}

% Vetor de rota (#1 = y, #2 = ids das cidades, #3/#4 = primeira/última posição invertida)
\tikzset{cid0/.style={bit,fill=white}, cid1/.style={bit,fill=red!25}}
\newcommand{\rotavec}[4]{%
	\foreach \v [count=\i] in {#2}{%
		\ifnum\i<#3 \def\st{cid0}\else\ifnum\i>#4 \def\st{cid0}\else\def\st{cid1}\fi\fi
		\node[\st] at (\i*0.65,#1) {\v};
	}%
}

\title[Aula 8 -- Metaheurísticas]{Metaheurísticas}
\subtitle{Escapando de ótimos locais}
\author[Santi]{Prof. \'Everton Santi}
\institute[UFRN/ECT]{ECT2524 -- Introdução à Otimização \\ Escola de Ciências e Tecnologia -- UFRN}
\date{Aula 8}

\begin{document}

\begin{frame}
	\titlepage
\end{frame}

\begin{frame}
	\frametitle{Sumário}
	\tableofcontents
\end{frame}

%------------------------------------------------------------
\section{Recapitulação}
%------------------------------------------------------------

\begin{frame}
	\frametitle{Heurística construtiva e busca local}
	\begin{itemize}
		\item \textbf{Heurística construtiva}: monta uma solução viável
			elemento a elemento (aleatória, gulosa, vizinho mais próximo);
		\item \textbf{Busca local}: parte de uma solução e a substitui
			por um vizinho melhor, enquanto houver;
		\item Vizinhança $\mathcal{N}(x)$: soluções obtidas de $x$ por uma
			pequena modificação;
		\item Exemplos: inversão de bit e troca (mochila); 2-opt (TSP).
	\end{itemize}
\end{frame}

\begin{frame}
	\frametitle{Ótimo local e ótimo global}
	\begin{columns}[T]
	\begin{column}{0.36\textwidth}
	\begin{itemize}
		\item Ótimo local: nenhum vizinho é melhor;
		\item Ótimo global: melhor solução de todo $S$;
		\item Um ótimo local depende da vizinhança adotada.
	\end{itemize}
	\end{column}
	\begin{column}{0.62\textwidth}
	\paisagemtres{%
			\addplot3[only marks,mark=*,black,mark size=2pt,
				nodes near coords={\scriptsize ótimo local},every node near coord/.style={above,black}]
				coordinates {(2,2,2.6)};
			\addplot3[only marks,mark=*,black,mark size=2pt,
				nodes near coords={\scriptsize ótimo local},every node near coord/.style={above,black}]
				coordinates {(8,2,2.6)};
			\addplot3[only marks,mark=*,black,mark size=2pt,
				nodes near coords={\scriptsize ótimo global},every node near coord/.style={above,black}]
				coordinates {(6,6,3.6)};
		}
	\end{column}
	\end{columns}
\end{frame}

\begin{frame}
	\frametitle{Busca local: só explotação}
	\begin{columns}[T]
	\begin{column}{0.36\textwidth}
	\begin{itemize}
		\item Sobe até o pico mais próximo;
		\item Outro ponto inicial, outro ótimo local.
	\end{itemize}
	\end{column}
	\begin{column}{0.62\textwidth}
	\paisagemtres{%
			\addplot3[only marks,mark=*,black,mark size=2pt,
				nodes near coords={\scriptsize $x_1$},every node near coord/.style={below,black}]
				coordinates {(0.9,1.0,0.6)};
			\addplot3[only marks,mark=*,black,mark size=2pt,
				nodes near coords={\scriptsize $x_2$},every node near coord/.style={below,black}]
				coordinates {(7.2,5.4,1.09)};
			\addplot3[only marks,mark=*,black,mark size=2pt,
				nodes near coords={\scriptsize $x_3$},every node near coord/.style={below,black}]
				coordinates {(8.9,1.0,0.78)};
			\addplot3[black,very thick,-{Stealth[length=2.5mm]}] coordinates {(0.9,1.0,0.6) (1.8,1.8,2.45)};
			\addplot3[black,very thick,-{Stealth[length=2.5mm]}] coordinates {(7.2,5.4,1.09) (6.2,5.9,3.4)};
			\addplot3[black,very thick,-{Stealth[length=2.5mm]}] coordinates {(8.9,1.0,0.78) (8.2,1.8,2.45)};
			\addplot3[only marks,mark=*,black,mark size=2pt,
				nodes near coords={\scriptsize $x_1'$},every node near coord/.style={above,black}]
				coordinates {(2,2,2.6)};
			\addplot3[only marks,mark=*,black,mark size=2pt,
				nodes near coords={\scriptsize $x_2'$},every node near coord/.style={above,black}]
				coordinates {(6,6,3.6)};
			\addplot3[only marks,mark=*,black,mark size=2pt,
				nodes near coords={\scriptsize $x_3'$},every node near coord/.style={above,black}]
				coordinates {(8,2,2.6)};
		}
	\end{column}
	\end{columns}
\end{frame}

\begin{frame}
	\frametitle{Limitação da busca local}
	\begin{itemize}
		\item Termina no primeiro ótimo local encontrado;
		\item Não dispõe de mecanismo para escapar dele;
		\item A qualidade do resultado depende da solução inicial;
		\item Nada garante que o ótimo local seja próximo do global.
	\end{itemize}
	\vspace{2mm}
	\begin{block}{Questão}
		Como continuar a busca a partir de um ótimo local, sem abandonar
		a busca local?
	\end{block}
\end{frame}

%------------------------------------------------------------
\section{Metaheurísticas}
%------------------------------------------------------------

\begin{frame}
	\frametitle{Definição de metaheurística}
	\begin{block}{Hillier e Lieberman (2010)}
		Método de resolução geral que fornece tanto uma estrutura quanto
		diretrizes de estratégia gerais para desenvolver um método
		heurístico específico que se ajuste a um tipo de problema
		particular.
	\end{block}
	\begin{block}{Hillier e Lieberman (2010)}
		Método de resolução geral que orquestra a interação entre
		procedimentos de melhoria local e estratégias de nível mais alto
		para criar um processo capaz de escapar dos ótimos locais e
		realizar uma busca consistente de uma região de soluções viáveis.
	\end{block}
\end{frame}

\begin{frame}
	\frametitle{Características}
	\begin{itemize}
		\item Gerenciam heurísticas de nível mais baixo e guiam a busca;
		\item Incorporam mecanismos para escapar de ótimos locais;
		\item Não dependem do problema: o que muda é a representação, a
			vizinhança e a função objetivo;
		\item Não fornecem prova de otimalidade.
	\end{itemize}
\end{frame}

\begin{frame}
	\frametitle{Elementos em comum}
	\begin{itemize}
		\item \textbf{Heurística construtiva}: gera soluções iniciais;
		\item \textbf{Busca local}: explotação, leva a um ótimo local;
		\item \textbf{Perturbação}: exploração, tira a busca do ótimo local;
		\item As metaheurísticas diferem na forma, na quantidade e na
			intensidade com que esses elementos são empregados.
	\end{itemize}
\end{frame}

\begin{frame}
	\frametitle{Perturbação}
	\begin{block}{Definição}
		Procedimento que gera uma nova solução $x'$, a partir de $x$, pela
		modificação aleatória de algumas de suas componentes.
	\end{block}
	\begin{itemize}
		\item A intensidade da modificação ($k$ componentes) é definida
			pelo usuário;
		\item Fuga de ótimos locais.
	\end{itemize}
	\vspace{2mm}
	\begin{center}
	\begin{tikzpicture}
		\node[anchor=east,font=\small] at (0.2,0.5) {solução $x$};
		\vetor{0.5}{1,0,1,0,1,1,0}
		\node[anchor=east,font=\small] at (0.2,-0.5) {perturbação, $k=2$};
		\vetor{-0.5}{0,0,1,0,0,1,0}
		\destaque{-0.5}{1}
		\destaque{-0.5}{5}
	\end{tikzpicture}
	\end{center}
\end{frame}

\begin{frame}
	\frametitle{Intensidade da perturbação}
	\begin{itemize}
		\item $k$ pequeno: $x'$ fica na região de atração de $x$ e a
			busca local retorna ao mesmo ótimo local;
		\item $k$ grande: $x'$ se aproxima de uma solução aleatória e a
			busca degenera em reinícios;
		\item Compromisso entre exploração e explotação;
		\item No TSP: reconexão de $k$ arestas (por exemplo, $k$ trocas de
			cidades).
	\end{itemize}
\end{frame}

%------------------------------------------------------------
\section{Busca local iterada (ILS)}
%------------------------------------------------------------

\begin{frame}
	\frametitle{Busca local iterada}
	\begin{itemize}
		\item \textit{Iterated Local Search} (ILS);
		\item Perturba o ótimo local atual e aplica nova busca local;
		\item Estrutura mais simples entre as metaheurísticas;
		\item Componentes: solução inicial, busca local, perturbação e
			critério de aceitação.
	\end{itemize}
\end{frame}

\begin{frame}
	\frametitle{Algoritmo ILS}
	\begin{enumerate}
		\item $x \leftarrow$ heurística construtiva;
		\item $x \leftarrow$ busca local$(x)$;
		\item Enquanto o critério de parada não for satisfeito:
		\begin{enumerate}
			\item $x' \leftarrow$ perturbação$(x)$;
			\item $x'' \leftarrow$ busca local$(x')$;
			\item Se $x''$ satisfaz o critério de aceitação: $x \leftarrow x''$;
			\item Atualiza a melhor solução encontrada;
		\end{enumerate}
		\item Retorna a melhor solução.
	\end{enumerate}
\end{frame}

\begin{frame}
	\frametitle{ILS na paisagem de soluções}
	\begin{center}
	\begin{tikzpicture}[xscale=0.85,yscale=0.65]
		\paisagem
		\fill (1.8,2.6) circle (2pt) node[above] {\scriptsize $x$};
		\fill (4.6,1.6) circle (2pt) node[below] {\scriptsize $x'$};
		\fill (5.8,3.6) circle (2pt) node[above] {\scriptsize $x''$};
		\draw[->,thick,red!70!black,bend left=35] (1.8,2.6) to (4.6,1.6);
		\draw[->,thick,black] (4.6,1.6) to (5.8,3.6);
		\node[font=\scriptsize,red!70!black] at (3.5,3.5) {perturbação};
		\node[font=\scriptsize] at (7.4,3.3) {busca local};
	\end{tikzpicture}
	\end{center}
\end{frame}

\begin{frame}
	\frametitle{Critério de aceitação}
	\begin{itemize}
		\item \textbf{Melhor}: aceita $x''$ somente se $f(x'') $ melhora $f(x)$;
			intensifica a busca em torno do ótimo local atual;
		\item \textbf{Sempre}: aceita qualquer $x''$; diversifica a busca,
			como um passeio aleatório entre ótimos locais;
		\item \textbf{Intermediário}: aceita $x''$ levemente pior, ou
			reinicia após certo número de iterações sem melhora.
	\end{itemize}
\end{frame}

\begin{frame}
	\frametitle{Exemplo: ILS na mochila}
	\begin{itemize}
		\item Instância da Aula 7 ($n = 7$, $C = 18$);
		\item Ótimo local da troca: $x$, benefício $83$.
	\end{itemize}
	\vspace{1mm}
	\begin{center}
	\begin{tikzpicture}
		\node[anchor=east,font=\small] at (0.2,0.9) {ótimo local $x$};
		\vetor{0.9}{1,0,1,0,1,1,0}
		\node[anchor=west,font=\small] at (5.0,0.9) {$83$};
		\node[anchor=east,font=\small] at (0.2,0.0) {perturbação $x'$};
		\vetor{0.0}{0,0,1,0,0,1,0}
		\destaque{0.0}{1}
		\destaque{0.0}{5}
		\node[anchor=west,font=\small] at (5.0,0.0) {$57$};
		\node[anchor=east,font=\small] at (0.2,-0.9) {busca local $x''$};
		\vetor{-0.9}{0,0,1,1,0,1,0}
		\destaque{-0.9}{4}
		\node[anchor=west,font=\small] at (5.0,-0.9) {$84$};
	\end{tikzpicture}
	\end{center}
	\begin{itemize}
		\item Remove os itens 1 e 5; a busca local insere o item 4;
		\item $x''$ é o ótimo global.
	\end{itemize}
\end{frame}

%------------------------------------------------------------
\section{GRASP}
%------------------------------------------------------------

\begin{frame}
	\frametitle{GRASP}
	\begin{itemize}
		\item \textit{Greedy Randomized Adaptive Search Procedure}
			(Feo e Resende, 1989);
		\item Repete: construção semigulosa e busca local;
		\item Mantém a melhor solução encontrada;
		\item A diversificação vem da construção, não de uma perturbação.
	\end{itemize}
\end{frame}

\begin{frame}
	\frametitle{Construção gulosa e semigulosa}
	\begin{itemize}
		\item \textbf{Gulosa}: a cada passo, insere o melhor candidato
			segundo o critério guloso;
		\item Determinística: sempre a mesma solução;
		\item \textbf{Semigulosa}: forma a \textbf{lista restrita de
			candidatos} (LRC) com os melhores candidatos e sorteia um
			deles;
		\item Cada construção produz uma solução diferente e de boa
			qualidade.
	\end{itemize}
\end{frame}

\begin{frame}
	\frametitle{Lista restrita de candidatos}
	\begin{itemize}
		\item $c_{\min}$, $c_{\max}$: menor e maior custo guloso entre os
			candidatos;
		\item Parâmetro $\alpha \in [0,1]$;
		\item Candidato $i$ entra na LRC se $c_i \leq c_{\min} +
			\alpha\,(c_{\max} - c_{\min})$;
		\item $\alpha = 0$: construção gulosa;
		\item $\alpha = 1$: construção aleatória.
	\end{itemize}
\end{frame}

\begin{frame}
	\frametitle{Guloso versus aleatório}
	\begin{center}
	\begin{tabular}{lcc}
		\toprule
		& $\alpha = 0$ (guloso) & $\alpha = 1$ (aleatório) \\
		\midrule
		Variância & nenhuma & alta \\
		Qualidade média & boa & baixa \\
		Convergência da busca local & rápida & lenta \\
		\bottomrule
	\end{tabular}
	\end{center}
	\begin{itemize}
		\item GRASP combina o melhor dos dois extremos.
	\end{itemize}
\end{frame}

\begin{frame}
	\frametitle{Algoritmo GRASP}
	\begin{enumerate}
		\item $f^* \leftarrow \infty$ (minimização);
		\item Enquanto o critério de parada não for satisfeito:
		\begin{enumerate}
			\item $x \leftarrow$ construção semigulosa$(\alpha)$;
			\item $x' \leftarrow$ busca local$(x)$;
			\item Se $f(x') < f^*$: $f^* \leftarrow f(x')$, $x^* \leftarrow x'$;
		\end{enumerate}
		\item Retorna $x^*$.
	\end{enumerate}
\end{frame}

%------------------------------------------------------------
\section{Simulated Annealing}
%------------------------------------------------------------

\begin{frame}
	\frametitle{Recozimento simulado}
	\begin{itemize}
		\item \textit{Simulated Annealing} (Kirkpatrick et al., 1983);
		\item Analogia com o recozimento de metais: aquecimento seguido de
			resfriamento lento;
		\item Temperatura alta: o sistema se move livremente;
		\item Temperatura baixa: o sistema se estabiliza em um estado de
			baixa energia;
		\item Energia $\leftrightarrow$ função objetivo.
	\end{itemize}
\end{frame}

\begin{frame}
	\frametitle{Aceitação de soluções piores}
	\begin{itemize}
		\item Sorteia um vizinho $x'$ de $x$ e calcula
			$\Delta = f(x') - f(x)$ (minimização);
		\item Se $\Delta \leq 0$: aceita $x'$;
		\item Se $\Delta > 0$: aceita $x'$ com probabilidade
			$e^{-\Delta/T}$ (critério de Metropolis);
		\item $T > 0$: temperatura;
		\item A perturbação é o próprio movimento aleatório de vizinhança.
	\end{itemize}
\end{frame}

\begin{frame}
	\frametitle{Probabilidade de aceitação}
	\begin{columns}[T]
	\begin{column}{0.38\textwidth}
	\begin{itemize}
		\item Piora grande: probabilidade menor;
		\item $T$ alto: aceita quase tudo (exploração);
		\item $T$ baixo: aceita quase só melhoras (explotação).
	\end{itemize}
	\end{column}
	\begin{column}{0.6\textwidth}
	\begin{tikzpicture}
	\begin{axis}[width=\linewidth,height=4.8cm,xlabel={$\Delta$},
		ylabel={$e^{-\Delta/T}$},xmin=0,xmax=5,ymin=0,ymax=1.05,
		legend pos=north east,legend style={font=\scriptsize},
		label style={font=\small},tick label style={font=\scriptsize}]
		\addplot[thick,red!70!black,domain=0:5,samples=40] {exp(-x/0.5)};
		\addplot[thick,blue!60!black,domain=0:5,samples=40] {exp(-x/1)};
		\addplot[thick,green!45!black,domain=0:5,samples=40] {exp(-x/3)};
		\legend{$T=0{,}5$,$T=1$,$T=3$}
	\end{axis}
	\end{tikzpicture}
	\end{column}
	\end{columns}
\end{frame}

\begin{frame}
	\frametitle{Resfriamento}
	\begin{itemize}
		\item Temperatura inicial $T_0$ alta;
		\item Esquema geométrico: $T \leftarrow \beta\,T$, com
			$\beta \in (0,1)$, por exemplo $0{,}95$;
		\item Em cada temperatura, um número fixo de iterações;
		\item Parada: $T$ abaixo de um limite ou número de iterações sem
			melhora.
	\end{itemize}
\end{frame}

\begin{frame}
	\frametitle{Algoritmo Simulated Annealing}
	\begin{enumerate}
		\item $x \leftarrow$ solução inicial; $T \leftarrow T_0$;
		\item Enquanto o critério de parada não for satisfeito:
		\begin{enumerate}
			\item Repete $L$ vezes:
			\begin{enumerate}
				\item Sorteia $x' \in \mathcal{N}(x)$; $\Delta \leftarrow f(x') - f(x)$;
				\item Se $\Delta \leq 0$ ou $\text{rand}(0,1) < e^{-\Delta/T}$: $x \leftarrow x'$;
				\item Atualiza a melhor solução encontrada;
			\end{enumerate}
			\item $T \leftarrow \beta\,T$;
		\end{enumerate}
		\item Retorna a melhor solução.
	\end{enumerate}
\end{frame}

%------------------------------------------------------------
\section{Algoritmo genético}
%------------------------------------------------------------

\begin{frame}
	\frametitle{Algoritmo genético}
	\begin{itemize}
		\item Inspirado na evolução biológica: seleção natural,
			reprodução, hereditariedade e mutação;
		\item Mantém uma \textbf{população} de soluções, não uma única;
		\item \textbf{Cromossomo}: representação de uma solução (por
			exemplo, vetor binário);
		\item \textbf{Aptidão} (\textit{fitness}): qualidade da solução,
			dada pela função objetivo;
		\item \textbf{Geração}: população em um dado instante.
	\end{itemize}
\end{frame}

\begin{frame}
	\frametitle{Operadores}
	\begin{itemize}
		\item \textbf{Seleção}: escolhe os pais; quanto maior a aptidão,
			maior a chance;
		\item \textbf{Cruzamento}: combina trechos de dois pais para
			gerar filhos;
		\item \textbf{Mutação}: altera aleatoriamente alguns genes do
			filho; corresponde à perturbação;
		\item Cruzamento: explotação das boas características dos pais;
			mutação: exploração.
	\end{itemize}
\end{frame}

\begin{frame}
	\frametitle{Cruzamento de um ponto}
	\begin{center}
	\begin{tikzpicture}
		\node[anchor=east,font=\small] at (0.2,1.4) {pai A};
		\vetor{1.4}{1,0,1,1,1,0,0}
		\node[anchor=east,font=\small] at (0.2,0.7) {pai B};
		\vetor{0.7}{0,0,1,0,1,1,0}
		\draw[red,very thick,dashed] (2.275,0.2) -- (2.275,1.9);
		\node[anchor=east,font=\small] at (0.2,-0.2) {filho 1};
		\vetor{-0.2}{1,0,1,0,1,1,0}
		\node[anchor=east,font=\small] at (0.2,-0.9) {filho 2};
		\vetor{-0.9}{0,0,1,1,1,0,0}
	\end{tikzpicture}
	\end{center}
	\begin{itemize}
		\item Ponto de corte sorteado;
		\item Filho 1: início de A e fim de B; filho 2: o inverso;
		\item Pode gerar solução inviável: reparo ou penalização.
	\end{itemize}
\end{frame}

\begin{frame}
	\frametitle{Algoritmo genético simples}
	\begin{enumerate}
		\item Gera a população inicial de $N$ cromossomos;
		\item Calcula a aptidão de cada cromossomo;
		\item Enquanto o critério de parada não for satisfeito:
		\begin{enumerate}
			\item Repete até formar $N$ filhos: seleciona dois pais,
				cruza com probabilidade $p_c$, aplica mutação com
				probabilidade $p_m$;
			\item Substitui a população pela nova geração;
			\item Calcula as aptidões e atualiza a melhor solução;
		\end{enumerate}
		\item Retorna a melhor solução.
	\end{enumerate}
\end{frame}

\begin{frame}
	\frametitle{Representação em outros problemas}
	\begin{itemize}
		\item Mochila: vetor binário;
		\item TSP: permutação das cidades; o cruzamento de um ponto
			gera rotas inválidas, exigindo operadores específicos;
		\item Variáveis contínuas: vetor de números reais.
	\end{itemize}
\end{frame}

%------------------------------------------------------------
\section{Outras metaheurísticas}
%------------------------------------------------------------

\begin{frame}
	\frametitle{Busca em vizinhança variável (VNS)}
	\begin{itemize}
		\item \textit{Variable Neighborhood Search} (Mladenović e
			Hansen, 1997);
		\item Família de vizinhanças $\mathcal{N}_1, \ldots,
			\mathcal{N}_{k_{\max}}$, de tamanho crescente;
		\item Perturbação: sorteia $x'$ em $\mathcal{N}_k(x)$, depois
			aplica busca local;
		\item Se melhora: $x \leftarrow x''$ e $k \leftarrow 1$; senão,
			$k \leftarrow k+1$;
		\item A intensidade da perturbação cresce até a busca sair do
			ótimo local;
		\item Variante VND: só a descida, sem perturbação.
	\end{itemize}
\end{frame}

\begin{frame}
	\frametitle{Busca tabu}
	\begin{itemize}
		\item Move-se sempre para o melhor vizinho, mesmo que pior que
			a solução atual;
		\item Lista tabu: movimentos recentes ficam proibidos por algumas
			iterações;
		\item Evita retornar ao ótimo local de que acabou de sair.
	\end{itemize}
\end{frame}

\begin{frame}
	\frametitle{Comparação}
	\begin{center}
	\small
	\begin{tabular}{lll}
		\toprule
		Método & Solução(ões) & Mecanismo de escape \\
		\midrule
		ILS & uma & perturbação + aceitação \\
		GRASP & várias independentes & construção semigulosa \\
		Simulated Annealing & uma & aceita piores, $e^{-\Delta/T}$ \\
		Algoritmo genético & população & cruzamento + mutação \\
		VNS & uma & vizinhança crescente \\
		\bottomrule
	\end{tabular}
	\end{center}
\end{frame}

%------------------------------------------------------------
\section{A tarefa (avaliativa)}
%------------------------------------------------------------

\begin{frame}
	\frametitle{Seminário (SEM1)}
	\begin{itemize}
		\item \textbf{25\% da nota}; equipes de até 3 pessoas (mesmas
			equipes das tarefas anteriores);
		\item Seminário em \textbf{23/10}, 15 minutos por equipe;
		\item Artigo dos anais do SBPO (qualquer edição) que aplique
			metaheurística a um problema;
		\item Apresentar como se fossem os autores no congresso;
		\item Envio do PDF na tarefa \textbf{SEM1} do SIGAA, até
			23/10, às 23h59.
	\end{itemize}
\end{frame}

\begin{frame}
	\frametitle{Conteúdo da apresentação}
	\begin{itemize}
		\item Problema, motivação e aplicação;
		\item Método de solução (metaheurística, representação,
			vizinhança ou operadores);
		\item Dados utilizados;
		\item Tempo de execução e máquina;
		\item Soluções ótimas ou não;
		\item Contribuições e limitações.
	\end{itemize}
	\vspace{2mm}
	\begin{alertblock}{Presença}
		Não apresentar, ou não assistir às apresentações dos colegas,
		implica não receber a nota do seminário.
	\end{alertblock}
\end{frame}

%------------------------------------------------------------
\section{Referências}
%------------------------------------------------------------

\begin{frame}
	\frametitle{Referências}
	\begin{itemize}
		\item Hillier, F. S.; Lieberman, G. J. \textit{Introduction to
			Operations Research}. 2010.
		\item Goldbarg, M. C.; Luna, H. P. L. \textit{Otimização
			Combinatória e Programação Linear}. 2005.
		\item Feo, T. A.; Resende, M. G. C. A probabilistic heuristic for
			a computationally difficult set covering problem.
			\textit{Operations Research Letters}, 8(2), 1989.
		\item Kirkpatrick, S.; Gelatt, C. D.; Vecchi, M. P. Optimization
			by simulated annealing. \textit{Science}, 220(4598), 1983.
		\item Mladenovi\'c, N.; Hansen, P. Variable neighborhood search.
			\textit{Computers \& Operations Research}, 24(11), 1997.
		\item Mitchell, M. \textit{An Introduction to Genetic
			Algorithms}. 1998.
	\end{itemize}
\end{frame}

\end{document}
